\documentclass[11pt]{article}
\usepackage[utf8]{inputenc}
\usepackage[T1]{fontenc}
\usepackage{geometry}
\usepackage{enumitem}
\usepackage{xcolor}
\usepackage{amsmath,amssymb,amsthm}
\usepackage{booktabs}
\usepackage{hyperref}

\geometry{margin=1in}

% ---------------------------------------------------------------------------
% Environments
% ---------------------------------------------------------------------------
\theoremstyle{definition}
\newtheorem{definition}{Definition}[section]
\newtheorem{assumption}{Assumption}
\renewcommand{\theassumption}{A\arabic{assumption}}
\newtheorem{example}[definition]{Example}
\theoremstyle{plain}
\newtheorem{theorem}[definition]{Theorem}
\newtheorem{proposition}[definition]{Proposition}
\theoremstyle{remark}
\newtheorem{remark}[definition]{Remark}

% Claims cited from memory and not yet checked against the source.
\newcommand{\verify}[1]{\textcolor{red}{[verify: #1]}}

% ---------------------------------------------------------------------------
% Notation
% ---------------------------------------------------------------------------
\newcommand{\X}{\mathcal{X}}
\newcommand{\Y}{\mathcal{Y}}
\newcommand{\Z}{\mathcal{Z}}
\newcommand{\W}{W}
\newcommand{\Prob}{\Delta}
\newcommand{\E}{\mathbb{E}}
\newcommand{\KL}{\mathrm{KL}}
\newcommand{\TV}{\mathrm{TV}}
\newcommand{\R}{\mathbb{R}}

\title{\textbf{The Mathematical Structure of Learning Problems}\\[6pt]
\large From records to the population loss}
\author{Working draft}
\date{\today}

\begin{document}
\maketitle

\noindent\emph{Part I draft. Every statement is a definition, a labelled
assumption, a cited theorem, or an example. Claims marked
\verify{\dots} are cited from memory and must be checked against the
source; see \texttt{verification.md}.}

% ============================================================================
\part{Framing}
% ============================================================================

% ============================================================================
\section{The problem}\label{sec:problem}
% ============================================================================

A learning problem starts with something we want a machine to do, and
a way of scoring how well it is done. Both exist before any data is
collected.

\begin{definition}[Decision problem]
A \emph{decision problem} consists of a set $A$ of actions, a set $\Y$
of outcomes, and a bounded utility $u(a,y)$ for taking action $a$ when
the outcome is $y$. The action is chosen after seeing an input $x$.
\end{definition}

\begin{example}
Classifying an image: $A=\Y$ is the set of labels and $u(a,y)=1$ if
$a=y$, else $0$. Screening for a tumour: $u$ penalises a missed tumour
more than a false alarm. Driving: $a$ is a steering command and $u$
scores the resulting trajectory.
\end{example}

Problems differ in how directly the utility is known.
\begin{itemize}[nosep]
\item \emph{Acting}: the actions and their consequences are specified
  (steering, treating a patient).
\item \emph{Predicting with a specified utility}: the action is a guess
  of the outcome and the utility scores being right, with possibly
  unequal costs for different errors (the tumour example).
\item \emph{Predicting with the utility left open}: in science one
  wants to know how the outcome is distributed because the uses come
  later; understanding cancer is valued for the drugs it will help to
  develop. Section~\ref{sec:utility} shows that this case leads to the
  KL loss.
\end{itemize}
In every case the learner is judged by the regret of its wrong
predictions or actions. The two cultures of statistical modelling
\cite{breiman2001}, parametric models whose parameters have scientific
meaning and flexible algorithmic models, solve the same problem; they
differ in the choice of model (Section~\ref{sec:models}). Two
refinements: when a parameter is itself the answer (an effect size used
to choose a dose), it must be identifiable, which is where the
structure studied in Part~II matters; and predicting the effect of an
intervention is regret under a different distribution from the one
that generated the records \cite{pearl2009}, which is outside the
scope of this note.

The utility is part of the problem and fixes which outcome matters. It
does not fix which inputs to record. That depends on what is already
known about the phenomenon: with little knowledge, one may record
inputs that carry almost no information about the outcome, and the
regularity of the records (Section~\ref{sec:regularity}) will then be
small or unrelated to the phenomenon of interest. When something is
known, experiments can be designed to maximise the expected information
the records carry \cite{lindley1956}.

% ============================================================================
\section{Records, and probability as the language}\label{sec:records}
% ============================================================================

What the machine receives are not the quantities of the world but
\emph{records} of them.

\begin{definition}[Observable, measurement, record]
An \emph{observable} $Z$ takes values in a space $\Z$ (often a
continuum, e.g.\ $\R^m$). A \emph{measurement system} is a Markov kernel
$\kappa(\,\cdot \mid z)$ from $\Z$ to a set of readings: it returns a
(possibly noisy) reading of $z$. An \emph{encoding} $e$ maps readings to
values in a set $\X$. A \emph{record} is $X = e(\text{reading})$. Because
every recorded value has finite precision, $\X$ is a finite set.
\end{definition}

Two facts about records are not assumptions: measurement devices
return different readings for the same observable (noise), and the
world varies in ways the records do not capture (hidden variables), so
equal inputs can be recorded with different outputs. Something must
describe this variability, and we choose probability.

\begin{remark}[Probability is a choice]
Probability is one of several mathematical languages for uncertainty,
and the most developed one. Arguments that it is forced by
coherence or by axioms of plausible reasoning are de Finetti's
coherence argument \cite{kadane2011} and Cox's theorem
\cite{vanhorn2003}; alternatives (imprecise probabilities, belief
functions, possibility measures, plausibility measures) and the
objections to these arguments are surveyed in \cite{halpern2017}. We
adopt probability for its mathematical development, not because the
alternatives are refuted.
\end{remark}

\begin{assumption}[Probabilistic representation]\label{as:prob}
There is a probability space $(\Omega,\mathcal{F},P)$ on which every
observable is a random variable, and measurement and encoding are
measurable. The records $X_1, X_2, \dots$ are then random variables
with values in $\X$.
\end{assumption}

The law of the records needs no further assumption.

\begin{proposition}[The law of the records exists]
Under Assumption~\ref{as:prob}, the record sequence has a law
$q_{1:\infty} = P\circ(X_1,X_2,\dots)^{-1}$ on $\X^{\mathbb N}$, and
every finite-sample law $q(x_{1:n})$ is one of its marginals.
Equivalently, any consistent family of finite-dimensional laws on the
finite set $\X$ defines a unique such law (Kolmogorov extension
\cite{kallenberg2002}).
\end{proposition}

\begin{remark}[Existence is free; stability is not]
Some probability law fits any data whatsoever (a point mass on the
observed sequence is one). Existence of $q_{1:\infty}$ therefore says
nothing about the world. The substantive question is whether the
\emph{same} law governs records not yet seen. That is what the
sampling symmetry of Section~\ref{sec:sampling} addresses.
\end{remark}

\begin{remark}[Time and actions]
The records generate a natural filtration
$\mathcal F_t = \sigma(X_1,\dots,X_t)$. A prediction of $X_{t+1}$ must
be $\mathcal F_t$-measurable (it may use only the past). If the learner
acts, the actions are recorded as part of $X_t$ and must also depend
only on past records; otherwise inference about the environment from
the records is not valid \cite{rubin1976}. For an unordered batch with
no actions the filtration imposes nothing.
\end{remark}

% ============================================================================
\section{Sampling symmetry}\label{sec:sampling}
% ============================================================================

\begin{definition}[Exchangeability]
Records $X_1,\dots,X_n$ are \emph{exchangeable} if their law is
unchanged by every permutation $s$ of $\{1,\dots,n\}$:
$q(x_1,\dots,x_n) = q(x_{s(1)},\dots,x_{s(n)})$. An infinite sequence
is exchangeable if every finite initial segment is. For supervised
records $X_i=(\mathrm{input}_i,\mathrm{output}_i)$ the pairs are
permuted together.
\end{definition}

\begin{theorem}[de Finetti; Hewitt--Savage \cite{hewitt1955,kallenberg2005}]
\label{thm:definetti}
Let $X_1, X_2,\dots$ be an infinite exchangeable sequence in a
standard Borel space $\X$. There is a unique probability measure $\mu$
on $\Prob(\X)$ such that for every $n$
\[
q(X_1\in A_1,\dots,X_n\in A_n) = \int_{\Prob(\X)} \prod_{i=1}^n Q(A_i)\,\mu(dQ).
\]
Moreover the empirical measures $\frac1n\sum_{i\le n}\delta_{X_i}$
converge almost surely to a random measure $Q$ with law $\mu$, and
conditionally on $Q$ the $X_i$ are i.i.d.\ with law $Q$.
\end{theorem}

\begin{definition}[The population distribution]
The realised value of $Q$ is written $q$ and called the
\emph{population distribution}. Conditionally on it, the records are
i.i.d.\ from $q$. The mixing measure $\mu$ is a prior on $\Prob(\X)$.
\end{definition}

So the textbook opening ``$X_1,\dots,X_n$ i.i.d.\ from $q$'' is a
\emph{consequence}: of Assumption~\ref{as:prob} and of exchangeability,
a symmetry stated about records.

\begin{theorem}[Finite exchangeability \cite{diaconis1980finite}]
If $X_1,\dots,X_n$ are exchangeable, the law of $(X_1,\dots,X_k)$ is
within total variation $2ck/n$ of a mixture of i.i.d.\ laws when $\X$
has $c$ elements, and within $k(k-1)/n$ in general.
\verify{norm convention for total variation}
\end{theorem}

A finite exchangeable sequence need not be a mixture of i.i.d.\ laws
(draw two balls without replacement from an urn holding one red and
one blue). The infinite theorem is what is used when the sample size
is sent to infinity; the finite one measures the error when the
population is finite.

\begin{remark}[Exchangeability can be designed and tested]
Random sampling, randomised collection order and shuffling make the
recorded sequence exchangeable by construction. Exchangeability can be
tested, though only rejected, never confirmed; conformal test
martingales give tests that are valid at every sample size
\cite{vovk2003testing}. Two different claims must not be confused:
exchangeability \emph{within} a dataset (which shuffling guarantees)
and exchangeability \emph{between} training and deployment records
(which it does not).
\end{remark}

\begin{remark}[Sequences]
When a record is a whole path (a document, an episode), exchangeability
is a statement about paths and Theorem~\ref{thm:definetti} gives a law
on path space. When there is a single long path, the corresponding
symmetry is \emph{stationarity} (invariance under time shift); a
stationary process is a mixture of ergodic ones, and along one path of
an ergodic process time averages converge to expectations (Birkhoff)
\cite{kallenberg2002}. In both regimes one realisation suffices to
identify the law in the limit.
\end{remark}

% ============================================================================
\section{Regularity and noise}\label{sec:regularity}
% ============================================================================

With $q$ in hand we can say what in the records is learnable. Records
are finite-valued, so all entropies below are finite.

\begin{proposition}[Decomposition of output uncertainty \cite{cover2006}]
For a record $(X,Y)\sim q$,
\[
\underbrace{H(Y)}_{\text{total}} \;=\; \underbrace{I(X;Y)}_{\text{regularity}} \;+\; \underbrace{H(Y\mid X)}_{\text{noise}} .
\]
\end{proposition}

\begin{definition}[Regularity and noise]
$I(X;Y)$, the information the input carries about the output, is the
\emph{regularity} of the problem; $H(Y\mid X)$ is its \emph{noise}. At
one extreme $I(X;Y)=0$: the input is useless and the best predictor is
the marginal $q(y)$ (which is not uniform guessing). At the other
$H(Y\mid X)=0$: the output is a deterministic function of the input.
\end{definition}

The same split holds for a single long sequence. For a stationary
process with block entropy $H(X_{1:L})$,
\[
H(X_{1:L}) \approx E + h\,L \qquad (L\to\infty),
\]
where $h$ is the entropy rate (noise per step) and the \emph{excess
entropy} $E = I(\text{past};\text{future})$ is the regularity
\cite{crutchfield2003,bialek2001}.

\begin{remark}[Regularity is compressibility]
By the Kraft inequality, a distribution $p$ on $\X$ is a code, and
$-\log_2 p(x)$ is the length of $x$ in bits \cite{cover2006}. Predicting
well under log loss and compressing well are therefore the same
thing, and regularity is what a code can exploit. For a single finite
dataset the same split is made by Kolmogorov's structure function:
total description length equals the length of a model plus the length
of the data's index within the model, and the data are incompressible
given a sufficient model \cite{gacs2001,vereshchagin2004}. Shannon
entropy is the average of Kolmogorov complexity
\cite{grunwald2003}. These algorithmic notions are uncomputable and
defined up to additive constants; we cite them as the rigorous form of
the idea, not as tools.
\end{remark}

\begin{remark}[Regularity is relative]
Regularity is always relative to a split: input versus output, or
past versus future. There is no absolute amount of regularity in a
dataset independent of what is to be predicted from what.
\end{remark}

% ============================================================================
\section{World symmetries produce regularity}\label{sec:world}
% ============================================================================

Sampling symmetry is about the order of records. A second kind of
symmetry is about the world itself.

\begin{definition}[World symmetry]
Let a group $G$ act on inputs, and (through a representation $\rho$) on
outputs. The problem is \emph{invariant} if $q(y \mid g\cdot x) =
q(y\mid x)$ and \emph{equivariant} if $q(\rho(g)\, y \mid g\cdot x) =
q(y\mid x)$, for all $g\in G$.
\end{definition}

\begin{remark}
World symmetries are usually properties of the conditional $q(y\mid
x)$, not of the marginal $q(x)$: whether an image shows a cat does not
depend on its orientation, but recorded photographs are mostly upright.
A symmetry of observables survives into records only if measurement
and encoding are equivariant; a pixel grid keeps only the four
right-angle rotations of the plane.
\end{remark}

A \emph{maximal invariant} $M$ is a function on inputs that is constant
on each orbit $\{g\cdot x: g\in G\}$ and takes different values on
different orbits.

\begin{theorem}[Invariance through a maximal invariant \cite{bloemreddy2020}]
\label{thm:noise-outsourcing}
Let $G$ be compact and act measurably on standard Borel spaces. Then
$q(y\mid x)$ is $G$-invariant if and only if $Y = f(\eta, M(X))$
almost surely for some measurable $f$ and some noise $\eta$ uniform on
$[0,1]$ and independent of $X$. \verify{exact hypotheses and theorem
number}
\end{theorem}

\begin{proposition}[Regularity lives on the quotient]
If $q(y\mid x)$ is $G$-invariant, then $Y$ and $X$ are conditionally
independent given $M(X)$, and so
\[
I(X;Y) = I(M(X);Y).
\]
\end{proposition}
\noindent This is the chain rule for mutual information applied to
Theorem~\ref{thm:noise-outsourcing}: $I(X;Y) = I(M(X);Y) + I(X;Y\mid
M(X))$ and the second term is zero. All of the regularity is carried
by \emph{which orbit} the input lies in; \emph{where} in the orbit it
lies carries none.

This is the precise sense in which symmetry produces regularity. The
converse is false: a deterministic trend is regular without any group
behind it. Building a known symmetry into a model is therefore a way to
spend model capacity only on the part of the input that carries
regularity, and doing so provably helps learning in settings where it
has been analysed \cite{elesedy2021}. \verify{scope of the strict
generalisation benefit}

For exchangeable records the same structure is fully explicit. The
orbit of a sequence under permutations is its \emph{type class} (all
sequences with the same empirical distribution); given its type, an
exchangeable sequence is uniform on its type class, and the type class
has about $2^{nH}$ elements, $H$ being the entropy of the type
\cite{csiszar1998}. The empirical distribution is the regularity and
the order is noise.

% ============================================================================
\section{Models}\label{sec:models}
% ============================================================================

A model is a family of candidate conditional distributions, indexed by
finitely many numbers.

\begin{definition}[Statistical model]
A \emph{statistical model} is a set $\W\subset\R^d$ of parameters
together with a map $w \mapsto p(\,\cdot\mid x, w)$ assigning to each
$w$ a conditional distribution of outputs given inputs. In practice the
map is a composition: a \emph{parametrisation} $w\mapsto f_w$ (the
architecture) followed by an \emph{output map} from $f_w(x)$ to a
distribution (e.g.\ a softmax). The \emph{image} of the model is the set
of conditional distributions it reaches.
\end{definition}

\begin{definition}[Realisability, fibres, identifiability]
The model is \emph{realisable} if $q(\,\cdot\mid x)$ is in its image
(for $q$-almost every $x$). The \emph{fibre} of $w$ is the set of
parameters giving the same conditional distribution as $w$ for
$q$-almost every input. The model is \emph{identifiable} if every fibre
is a single point.
\end{definition}

Fibres depend on the input distribution: two parameters may agree on
every input that occurs and disagree elsewhere.

\begin{definition}[Regular model]
A model is \emph{regular} if it is identifiable, the map
$w\mapsto p(\,\cdot\mid x,w)$ is smooth, and the Fisher information
matrix
\[
I(w) = \E\big[\nabla_w \log p(Y\mid X,w)\,\nabla_w \log p(Y\mid X,w)^{\top}\big]
\]
is positive definite at every $w$ \cite{watanabe2009}. A model that is
not regular is \emph{singular}.
\end{definition}

\begin{example}
The Gaussian model $\W=\R\times(0,\infty)$, $(m,s)\mapsto N(m,s^2)$, is
regular. The model $(a,b,s)\mapsto N(ab,s^2)$ has the same image and is
singular: at $a=b=0$ the fibre is the cross $\{a=0\}\cup\{b=0\}$ and
the Fisher information is degenerate. Singularity is a property of the
map, not of the set of distributions it reaches.
\end{example}

A good model is one whose image comes close to $q$, that is, one that
captures the regularity of Section~\ref{sec:regularity}. To say how
close, we need to say what the model will be used for.

% ============================================================================
\section{Utility, and why the loss is KL}\label{sec:utility}
% ============================================================================

We return to the decision problem of Section~\ref{sec:problem}. Given
input $x$, the best action under a distribution $p(\,\cdot\mid x)$ is
$a_p(x) \in \arg\max_{a} \E_{p(y\mid x)}[u(a,y)]$.
The action we want is $a_q(x)$, which depends only on the utility and
on the unknown $q$. Acting with a model means acting with $a_p(x)$ in
its place, and the utility judges the model by the resulting loss of
expected utility (the \emph{regret}).

\begin{proposition}[KL controls every bounded utility]\label{prop:regret}
If $u$ takes values in an interval of length $R$, then
\[
\E_{q(y\mid x)}\big[u(a_q(x),y)\big] - \E_{q(y\mid x)}\big[u(a_p(x),y)\big]
\;\le\; R\,\sqrt{2\,\KL\big(q(\cdot\mid x)\,\|\,p(\cdot\mid x)\big)} .
\]
\end{proposition}
\noindent For any action, the expected utilities under $q$ and $p$
differ by at most $R\cdot\TV(q,p)$. Adding and subtracting expected
utilities under $p$, and using that $a_p$ is optimal under $p$, the
regret is at most $2R\cdot\TV$. Pinsker's inequality
$\TV\le\sqrt{\KL/2}$ \cite{cover2006} gives the bound. Averaging over
inputs and applying Jensen's inequality, the expected regret is at most
$R\sqrt{2\,\E_X \KL(q(\cdot\mid X)\,\|\,p(\cdot\mid X))}$.

So a model with small KL to $q$ is good for \emph{every} bounded
decision problem at once, including ones not yet specified. This is the
situation of predicting with the utility left open
(Section~\ref{sec:problem}), where the eventual uses are unknown. For
that case there is also a direct utility: state a distribution $p$ and
receive $\log p(y)$ when $y$ occurs (the logarithmic score). Its
expected regret is exactly $\KL(q\,\|\,p)$, by Gibbs' inequality, and
under smoothness and locality conditions it is the only utility for
stating distributions that rewards honesty; Bernardo proposed it as the
utility of pure scientific inference \cite{bernardo1979}. \verify{exact uniqueness
conditions in Bernardo 1979}

\begin{remark}[A fixed task may need less than $q$]
A particular utility often depends on only part of $q$. With the 0--1
utility of classification, $a_q(x)$ is the most probable label, so a
model whose most probable label always matches $q$'s is optimal for
that task however large its KL loss. This gives two strategies. A
\emph{task-specific} learner optimises the given utility directly and
models only what it needs (Vapnik's advice not to solve a harder
problem than the one posed). A \emph{general-purpose} learner models
all of $q$ under KL and, by Proposition~\ref{prop:regret}, serves every
bounded utility at once, including utilities not yet specified. KL is
the right loss when the task is not fixed in advance.
\end{remark}

\begin{definition}[Population loss]
For a model $p(\,\cdot\mid x,w)$, the \emph{population loss} and the
\emph{KL loss} are
\[
L(w) = \E_q\big[-\log p(Y\mid X,w)\big], \qquad
K(w) = \E_X\,\KL\big(q(\cdot\mid X)\,\|\,p(\cdot\mid X,w)\big).
\]
\end{definition}

\begin{proposition}[KL is the regularity not captured]
\[
L(w) = H(Y\mid X) + K(w), \qquad\text{equivalently}\qquad
H(Y) - L(w) = I(X;Y) - K(w).
\]
\end{proposition}
\noindent The first identity is the definition of KL; the second
subtracts it from $H(Y)=I(X;Y)+H(Y\mid X)$. The left side of the second
is how much better the model does than knowing nothing about the input
beyond the output frequencies. It equals the regularity minus the KL
loss. A model captures all the regularity exactly when $K(w)=0$, and
then its loss is the noise $H(Y\mid X)$, which no model can beat.

Neither $L$ nor $K$ can be computed, since both involve $q$. What can
be computed from records is the \emph{empirical loss}
\[
L_n(w) = -\frac1n \sum_{i=1}^n \log p(y_i\mid x_i,w),
\]
which converges to $L(w)$ for each $w$ by the strong law of large
numbers applied conditionally on $q$. Since $L = H(Y\mid X) + K$ and
$H(Y\mid X)$ does not depend on $w$, minimising $L$ and minimising $K$
are the same problem. For
sequences, the chain rule writes the log-likelihood of a path as a sum
of next-step conditional log-likelihoods, so the same formula holds
with each record a step given its history.

This is the standard setup of statistical learning: records i.i.d.\
from $q$, a model $w\mapsto p(\,\cdot\mid x,w)$, and the loss $K$
estimated by $L_n$. Part I has shown where each piece comes from.

\begin{table}[h]
\centering
\small
\begin{tabular}{@{}lll@{}}
\toprule
Ingredient & Status & Where \\
\midrule
task and utility $u$ & given & \S\ref{sec:problem} \\
measurement system, encoding & chosen & \S\ref{sec:records} \\
probability as the language & chosen & \S\ref{sec:records} \\
probability space (A1) & assumed & \S\ref{sec:records} \\
law of the records $q_{1:\infty}$ & follows from A1 & \S\ref{sec:records} \\
exchangeability / stationarity & assumed; designable, testable & \S\ref{sec:sampling} \\
population distribution $q$ & follows (de Finetti) & \S\ref{sec:sampling} \\
regularity $I(X;Y)$, noise $H(Y\mid X)$ & properties of $q$ & \S\ref{sec:regularity} \\
world symmetry $G$ & assumed (domain knowledge) & \S\ref{sec:world} \\
model $w\mapsto p(\cdot\mid x,w)$ & chosen & \S\ref{sec:models} \\
realisability & assumed or not & \S\ref{sec:models} \\
loss $K$ & follows from the utility & \S\ref{sec:utility} \\
$K$ replaced by $L_n$ & approximation (statistical) & \S\ref{sec:utility} \\
\bottomrule
\end{tabular}
\caption{Ledger for Part I.}
\end{table}

% ============================================================================
\begin{thebibliography}{99}
% ============================================================================

\bibitem{bernardo1979} J.~M.~Bernardo. Expected information as expected
utility. \emph{Annals of Statistics} 7(3), 1979.

\bibitem{bialek2001} W.~Bialek, I.~Nemenman, N.~Tishby. Predictability,
complexity, and learning. \emph{Neural Computation} 13, 2001.

\bibitem{bloemreddy2020} B.~Bloem-Reddy, Y.~W.~Teh. Probabilistic
symmetries and invariant neural networks. \emph{JMLR} 21, 2020.

\bibitem{breiman2001} L.~Breiman. Statistical modeling: the two cultures.
\emph{Statistical Science} 16(3), 2001.

\bibitem{cover2006} T.~M.~Cover, J.~A.~Thomas. \emph{Elements of
Information Theory}, 2nd ed. Wiley, 2006.

\bibitem{crutchfield2003} J.~P.~Crutchfield, D.~P.~Feldman. Regularities
unseen, randomness observed: levels of entropy convergence.
\emph{Chaos} 13, 2003.

\bibitem{csiszar1998} I.~Csisz\'ar. The method of types. \emph{IEEE
Trans.\ Information Theory} 44(6), 1998.

\bibitem{diaconis1980finite} P.~Diaconis, D.~Freedman. Finite
exchangeable sequences. \emph{Annals of Probability} 8(4), 1980.

\bibitem{elesedy2021} B.~Elesedy, S.~Zaidi. Provably strict
generalisation benefit for equivariant models. \emph{ICML}, 2021.

\bibitem{gacs2001} P.~G\'acs, J.~Tromp, P.~Vit\'anyi. Algorithmic
statistics. \emph{IEEE Trans.\ Information Theory} 47(6), 2001.

\bibitem{grunwald2003} P.~Gr\"unwald, P.~Vit\'anyi. Kolmogorov complexity
and information theory, with an interpretation in terms of questions
and answers. \emph{J.\ Logic, Language and Information} 12, 2003.

\bibitem{halpern2017} J.~Y.~Halpern. \emph{Reasoning About Uncertainty},
2nd ed. MIT Press, 2017.

\bibitem{hewitt1955} E.~Hewitt, L.~J.~Savage. Symmetric measures on
Cartesian products. \emph{Trans.\ AMS} 80, 1955.

\bibitem{kadane2011} J.~B.~Kadane. \emph{Principles of Uncertainty}.
Chapman \& Hall/CRC, 2011.

\bibitem{kallenberg2002} O.~Kallenberg. \emph{Foundations of Modern
Probability}, 2nd ed. Springer, 2002.

\bibitem{kallenberg2005} O.~Kallenberg. \emph{Probabilistic Symmetries
and Invariance Principles}. Springer, 2005.

\bibitem{lindley1956} D.~V.~Lindley. On a measure of the information
provided by an experiment. \emph{Annals of Mathematical Statistics}
27(4), 1956.

\bibitem{pearl2009} J.~Pearl. \emph{Causality}, 2nd ed. Cambridge
University Press, 2009.

\bibitem{rubin1976} D.~B.~Rubin. Inference and missing data.
\emph{Biometrika} 63(3), 1976.

\bibitem{vanhorn2003} K.~S.~Van Horn. Constructing a logic of plausible
inference: a guide to Cox's theorem. \emph{Int.\ J.\ Approximate
Reasoning} 34, 2003.

\bibitem{vereshchagin2004} N.~Vereshchagin, P.~Vit\'anyi. Kolmogorov's
structure functions and model selection. \emph{IEEE Trans.\
Information Theory} 50(12), 2004.

\bibitem{vovk2003testing} V.~Vovk, I.~Nouretdinov, A.~Gammerman. Testing
exchangeability on-line. \emph{ICML}, 2003.

\bibitem{watanabe2009} S.~Watanabe. \emph{Algebraic Geometry and
Statistical Learning Theory}. Cambridge University Press, 2009.

\end{thebibliography}

\end{document}
